\section{Motivation}

% OWNER: Kelly Fu (intro per proposal)

\begin{frame}{Motivation: Can grammar survive noisy input?}
  % TODO (Kelly): children acquire robust SVA despite noisy input (speech
  % errors, non-native speech, agreement violations).
  % TODO (Kelly): poverty-of-the-stimulus framing -> how robust are LMs?
\begin{itemize}
\item Children acquire robust grammatical systems despite hearing noisy and imperfect language
\item This observation motivates the classic \textbf{Poverty of the Stimulus} debate:

      Is grammar learned from data, or supported by innate structure?
\item Human experiments cannot systematically manipulate children's linguistic input
\item Neural language models provide a controllable testbed:
\begin{itemize}
\item architecture = innate structure
\item training corpus = linguistic environment
\end{itemize}

\end{itemize}

\end{frame}


\begin{frame}{Research Questions}

\textbf{RQ1: Can grammar be learned from noisy input?}

\begin{itemize}
    \item How much grammatical corruption can a model tolerate?
    \item Is there a critical noise threshold beyond which acquisition breaks down?
\end{itemize}


\textbf{RQ2: How does noise change generalization?}

\begin{itemize}
    \item Does the model maintain a hierarchical syntactic rule?
    \item Or does it fall back on a simpler linear nearest-noun heuristic?
\end{itemize}


\textbf{RQ3: Does structure buy robustness?}

\begin{itemize}
    \item Compare LSTM and RNNG under identical noise conditions.
    \item Does explicit syntactic bias help preserve grammatical learning?
\end{itemize}

\end{frame}
